\section{Results}\label{sec:results}

\subsection{Generation quality of prior models}\label{subsec:prior_generation_quality}

We first evaluated the basic generation quality of two PepALD prior models: PepALD$_{\mathrm{ChEMBL}}$, obtained by pretraining on ChEMBL32 HELM peptide sequences, and PepALD$_{\mathrm{CycPeptMPDB}}$, obtained by further fine-tuning PepALD$_{\mathrm{ChEMBL}}$ on CycPeptMPDB cyclic peptide sequences.
We compared these two PepALD priors with HELM-GPT and PepMDLM.
HELM-GPT provides a closely matched HELM-level baseline, as it was likewise pretrained on the ChEMBL database and further fine-tuned on the CycPeptMPDB database \citep{ref_helm_gpt}.
PepMDLM denotes the SMILES-based peptide diffusion prior underlying PepTune before multi-objective inference-time guidance; this prior was trained on a peptide SMILES corpus compiled from CycPeptMPDB, SmProt, and a CycloPs-generated modified-peptide library \citep{ref_peptune}.

Five generation-quality metrics are adapted from MOSES \citep{ref_moses} to evaluate generation quality: validity, uniqueness, novelty, internal diversity, and similarity to the nearest training-set neighbor (SNN). Detailed definitions and implementation details are provided in Supplementary Section~8. Table~\ref{tab:prior_models} presents the evaluation results of the PepALD prior models and reference baselines.
Benefiting from HELM-based monomer-level assembly, a chemically informed foundation-model prior, and autoregressive diffusion for chain-like molecular generation, PepALD$_{\mathrm{ChEMBL}}$ achieved a validity of 0.951 for generating chemically valid peptide structures, substantially higher than HELM-GPT$_{\mathrm{ChEMBL}}$ (validity of 0.708) and PepMDLM (validity of 0.334).
Beyond validity, benefiting from the generative capability of the diffusion-based paradigm, the ChEMBL-pretrained PepALD prior, together with PepMDLM, achieved peak performance in both novelty and uniqueness, with values of 1.000.
It also achieved the highest internal diversity of 0.790 among the baseline models, along with a SNN of 0.441, outperforming HELM-GPT prior model and PepMDLM.
After fine-tuning on the CycPeptMPDB dataset, PepALD continued to show leading performance across these generation-quality metrics.

\begin{table*}[!t]
\centering
\caption{Evaluation of PepALD prior models and reference baselines.}
\label{tab:prior_models}
\small
\setlength{\tabcolsep}{0pt}
\begin{tabular*}{\textwidth}{@{\extracolsep{\fill}}lccccc@{}}
\toprule
Model & Validity & Uniqueness & Diversity & SNN & Novelty \\
\midrule
Baseline$_{\mathrm{ChEMBL}}$ & \textbf{1.000} & \textbf{1.000} & \underline{0.768} & 1.000 & 0.000 \\
HELM-GPT$_{\mathrm{ChEMBL}}$ & 0.708 & 0.890 & 0.740 & 0.750 & \underline{0.889} \\
HELM-GPT$_{\mathrm{CycPeptMPDB}}$ & 0.839 & \underline{0.913} & 0.595 & 0.975 & 0.461 \\
PepMDLM & 0.334 & \textbf{1.000} & 0.709 & 0.590 & \textbf{1.000} \\
PepALD$_{\mathrm{ChEMBL}}$ & \underline{0.951} & \textbf{1.000} & \textbf{0.790} & \textbf{0.441} & \textbf{1.000} \\
PepALD$_{\mathrm{CycPeptMPDB}}$ & 0.933 & \textbf{1.000} & 0.726 & \underline{0.504} & \textbf{1.000} \\
\bottomrule
\end{tabular*}
\begin{flushleft}
\footnotesize Baseline$_{\mathrm{ChEMBL}}$ denotes samples from the ChEMBL dataset. PepALD$_{\mathrm{ChEMBL}}$ and PepALD$_{\mathrm{CycPeptMPDB}}$ are the PepALD prior models, corresponding to the ChEMBL-pretrained and CycPeptMPDB-fine-tuned generators, respectively. HELM-GPT$_{\mathrm{ChEMBL}}$ and HELM-GPT$_{\mathrm{CycPeptMPDB}}$ are the corresponding HELM-GPT prior models; PepMDLM denotes the pretrained PepTune/PepMDLM reference model.
\end{flushleft}
\end{table*}

\begin{figure*}[!t]
\centering
\includegraphics[width=\textwidth]{result_data/figure_about_monomer/fig2_prior_monomer_coverage.pdf}
\caption{Prior monomer-space and residue-family coverage before downstream optimization.
(a) Monomer-level chemical-space coverage, where the gray background denotes the full monomer library projected from molecular fingerprints and highlighted points denote monomers used by HELM-GPT and PepALD samples; point size is proportional to monomer usage frequency.
(b) Monomer rarefaction curves over generated peptides.
(c) Fingerprint-cluster rarefaction curves over generated peptides; fingerprint clusters were obtained by applying K-means clustering to SVD-reduced Morgan-fingerprint representations of the full mapped monomer library.
Only clean monomers mapped to the monomer SMILES library are included.}
\label{fig:prior_monomer_coverage}
\end{figure*}

Because PepALD performs diffusion-based generation in a latent space informed by molecular chemical properties, it is not constrained by local minima arising from fixed combinations of building blocks, enabling the generation of diverse peptide molecules.
We therefore examined whether its generalizability was accompanied by broader use of the underlying monomer library.
In the 1000-sample monomer-level evaluation, PepALD used 263 clean monomers mapped to the monomer SMILES library, whereas HELM-GPT used 34 mapped monomers (Fig.~\ref{fig:prior_monomer_coverage}a).
To assess whether this broader monomer usage also covered distinct residue families, we grouped the mapped monomer library into 48 fingerprint-derived clusters.
The rarefaction curves show that this advantage was sustained across sampling depths at both resolutions: PepALD accumulated new monomers more rapidly (Fig.~\ref{fig:prior_monomer_coverage}b) and occupied 37 clusters, compared with 15 clusters for HELM-GPT (Fig.~\ref{fig:prior_monomer_coverage}c).
Together, the monomer- and monomer-cluster-level analyses indicate that PepALD samples more broadly within the existing monomer chemical space, covering more residue identities and fingerprint-derived residue families.

\subsection{Cell permeability optimization}\label{subsec:permeability_optimization}

To establish a permeability-enriched starting point for subsequent receptor-specific optimization, we next evaluated whether PepALD could be shifted toward cell-permeable macrocyclic peptides while retaining chemical diversity.
As described in Section~\ref{subsubsec:permeability_prior}, the permeability-enriched PepALD model was initialized from the cyclic prior and further fine-tuned on the top $1\,000$ training-set molecules with the highest predicted permeability scores.
For permeability scoring, we adopted the random-forest predictor previously used in HELM-GPT \citep{ref_helm_gpt}, which achieved a Spearman correlation of 0.82 on a held-out macrocyclic-peptide test set and successfully guided HELM-space permeability optimization.
Accordingly, we used this predictor to compute the average predicted permeability (less negative values indicating better membrane permeability) of cyclic peptides generated by each method.
During cyclic-prior training followed by permeability-oriented fine-tuning, the fitted mean permeability increased from $-7.069$ to $-6.688$ and then to $-5.822$ (Supplementary Fig.~S2), indicating a progressive shift toward the high-permeability regime.

\begin{table*}[!b]
\centering
\small
\caption{Evaluation of permeability-oriented generation.}
\label{tab:permeability_models}
\setlength{\tabcolsep}{0pt}
\begin{tabular*}{\textwidth}{@{\extracolsep{\fill}}lccccc@{}}
\toprule
Model & Permeability & Solubility & Diversity & SNN & Validity \\
\midrule
Baseline$_{\mathrm{ChEMBL}}$ & -7.104 & \underline{0.837} & \underline{0.768} & 1.000 & \textbf{1.000} \\
HELM-GPT$_{\mathrm{perm}}$ & \textbf{-5.396} & 0.759 & 0.553 & 0.621 & 0.253 \\
PepTune & -7.378 & 0.825 & 0.611 & 0.537 & 0.341 \\
PepALD$_{\mathrm{perm}}$ & \underline{-5.822} & 0.788 & 0.660 & 0.585 & \underline{0.955} \\
PepALD$_{\mathrm{ChEMBL}}$ & -7.069 & 0.821 & \textbf{0.790} & \textbf{0.441} & 0.951 \\
PepALD$_{\mathrm{CycPeptMPDB}}$ & -6.669 & \textbf{0.840} & 0.726 & \underline{0.504} & 0.933 \\
\bottomrule
\end{tabular*}
\begin{flushleft}
\footnotesize Baseline$_{\mathrm{ChEMBL}}$ denotes samples from the ChEMBL dataset. PepALD$_{\mathrm{ChEMBL}}$ and PepALD$_{\mathrm{CycPeptMPDB}}$ are the PepALD priors before permeability fine-tuning. HELM-GPT$_{\mathrm{perm}}$ is the permeability-oriented HELM-GPT baseline; PepTune denotes the PepTune/PepMDLM reference model. Permeability is higher-is-better, and solubility is the PepTune predicted soluble-class probability.
\end{flushleft}
\end{table*}

Table~\ref{tab:permeability_models} compares PepALD$_{\mathrm{perm}}$ with permeability-oriented baselines.
HELM-GPT$_{\mathrm{perm}}$ denotes the HELM-GPT model optimized for permeability, whereas PepTune represents the PepMDLM model using permeability-oriented guidance.
At the peptide level, the permeability-oriented PepALD samples occupied a broad region while retaining partial overlap with HELM-GPT (Fig.~\ref{fig:permeability_peptide_space}), indicating that the permeability shift did not collapse generation into a narrow peptide-space neighborhood.

\begin{figure}[!t]
\centering
\includegraphics[width=0.70\linewidth]{result_data/figure_about_monomer/fig3_peptide_space.pdf}
\caption{Peptide-level chemical-space coverage after permeability-oriented generation, shown as an MDS projection of pairwise molecular-fingerprint distances for valid samples from PepMDLM, HELM-GPT, and PepALD.}
\label{fig:permeability_peptide_space}
\end{figure}

Because adequate aqueous solubility is needed to maintain the concentration gradient that drives membrane flux, we also report the predicted soluble-class probability alongside permeability \citep{ref_peptune}, with the valid-sample solubility distributions shown in Fig.~\ref{fig:permeability_optimization}a.
In terms of permeability, PepALD$_{\mathrm{perm}}$ achieved a mean predicted permeability of -5.8215, marginally inferior to HELM-GPT$_{\mathrm{perm}}$ (-5.3964) among the generative models and substantially better than PepTune (-7.3783), as illustrated by the predicted permeability distributions in Fig.~\ref{fig:permeability_optimization}b.
However, the higher permeability of HELM-GPT$_{\mathrm{perm}}$ was accompanied by a marked reduction in generation usability.
Its validity was only 0.253, compared with 0.955 for PepALD$_{\mathrm{perm}}$ and 0.341 for PepTune.
PepALD$_{\mathrm{perm}}$ also retained stronger generation-quality characteristics, with an internal diversity of 0.660 that exceeded both HELM-GPT$_{\mathrm{perm}}$ and PepTune, together with a moderate SNN of 0.585.

The importance of validity and diversity was further supported by the SNN-gated batch-discovery analysis in Fig.~\ref{fig:permeability_optimization}c.
When equal-size raw-sample batches were filtered by an SNN $<0.60$ novelty gate, HELM-GPT$_{\mathrm{perm}}$ produced progressively fewer retained candidates during repeated sampling, suggesting that its generated samples increasingly revisited chemically similar neighborhoods.
In contrast, PepALD$_{\mathrm{perm}}$ continued to provide nonredundant candidates within the desired permeability range across sampling batches.

Through a permeability-oriented strategy, PepALD$_{\mathrm{perm}}$ not only exhibits superior properties in terms of permeability scores, but also preserves the structural plausibility, diversity, and solubility of the generated cyclic peptides, which suggests that PepALD$_{\mathrm{perm}}$ provides an effective and balanced framework for the in silico design of cyclic peptides.

\begin{figure*}[!t]
\centering
\includegraphics[width=0.97\textwidth]{result_data/figure_solu_perm/fig4_abc_property_panels.pdf}
\caption{Permeability-oriented optimization and sample usability.
(a) Predicted soluble-class probability distributions for valid samples from PepTune, HELM-GPT$_{\mathrm{perm}}$, and PepALD$_{\mathrm{perm}}$; non-ok conversions or NaN scores were excluded.
(b) Predicted permeability-score distributions of valid generated samples from the same generators; sentinel scores ($\leq -10$) were excluded.
(c) SNN-gated batch discovery from equal-size raw-sample batches.
Failed HELM-to-molecule conversions were discarded.
Each batch drew 40 samples per model, and the remaining candidates were retained only if their maximum Morgan-fingerprint Tanimoto similarity to the previously accepted pool was below 0.60.
For readability, panel (c) displays retained candidates with permeability scores between $-7$ and $-5$.}
\label{fig:permeability_optimization}
\end{figure*}

\subsection{Optimization for specific tasks}\label{subsec:specific_task_optimization}

We further applied \mytool to two target-specific macrocyclic peptide optimization cases: SPRY domain-containing SOCS box protein 2 (SPSB2, PDB id: 6DN5) and Chorismate Mutase from Mycobacterium tuberculosis (MtbCM, PDB id: 9BT3).
To optimize and evaluate the ability of generated macrocyclic peptides to bind the target proteins, we utilized docking scores from Uni-Dock, a GPU-accelerated Vina-style docking engine, as both the candidate rewards and the affinity metric \citep{ref_unidock}.
To more comprehensively evaluate the applicability of PepALD in macrocyclic peptide drug-design scenarios, the predicted membrane permeability and aqueous solubility of the generated peptides were also incorporated as evaluation metrics.

\subsubsection{Cyclic peptide generation for SPSB2}\label{subsubsec:spsb2_generation}

SPSB2 is a SPRY domain-containing adaptor protein that recruits iNOS for ubiquitin-mediated degradation, thereby suppressing nitric oxide production and modulating host antimicrobial responses.
Designing cyclic peptide inhibitors against the SPSB2-iNOS interface can potently and selectively stabilize iNOS levels, offering a structure-guided strategy for enhancing nitric oxide-dependent clearance of intracellular pathogens.

In this task, we aim to generate cell-permeable macrocyclic peptide binders targeting SPSB2 structure from PDB structure 6DN5 using \mytool.
Given the exploratory capacity of \mytool, we did not introduce any additional 6DN5-specific peptide or ligand set as prior information.
Instead, the WP-DPO policy was initialized directly from PepALD$_{\mathrm{perm}}$ and optimized using the strategy described in Section~\ref{subsec:preference_optimization}.

The Vina-score distributions in Fig.~\ref{fig:specific_task_optimization}a show a steady shift toward more favorable docking scores during WP-DPO.
Through ten optimization rounds, \mytool can generate macrocyclic peptide with more favorable docking scores, improving the mean Vina score from -5.832 for PepALD$_{\mathrm{perm}}$ to -7.020 for PepALD$_{\mathrm{SPSB2}}$.
The top-ranked subsets exceeded the reference ligand in predicted docking affinity, with mean Vina scores of -8.037 for the top $10\%$ and -7.746 for the top $25\%$, compared with -7.627 for the reference ligand.
At the same time, the optimized population retained high predicted solubility (0.863) and internal diversity (0.654), and its predicted permeability (-6.593) remained higher than that of the reference ligand (-7.181).

\subsubsection{Cyclic peptide generation for MtbCM}\label{subsubsec:mtbcm_generation}

Chorismate mutase from Mycobacterium tuberculosis (Mtb) is a secreted enzyme implicated in bacterial virulence and host-pathogen interactions. 
Designing cyclic peptide inhibitors against Mtb chorismate mutase (MtbCM) has the potential of enabling potent enzymatic inhibition and providing new ligand scaffolds for tuberculosis therapies amid rising antibiotic resistance.

In this task, we used \mytool to generate cyclic peptides against the MtbCM structure represented by PDB id 9BT3. The WP-DPO optimization procedure followed the same workflow as the SPSB2-binding cyclic peptide generation task described above.

The MtbCM-targeted task showed the same self-improving behavior under generated-data WP-DPO.
The mean Vina score improved from \mbox{-7.641} for PepALD$_{\mathrm{perm}}$ to \mbox{-8.645} for PepALD$_{\mathrm{MtbCM}}$, and the top-ranked generated subsets reached substantially stronger docking scores.
In particular, the top $10\%$ subset achieved a mean Vina score of -10.324, and the top $25\%$ subset reached -9.781, matching the reference ligand-level docking performance (-9.730).
The optimized cyclic peptides exhibited substantially better predicted permeability than the reference ligand, with a mean score of -5.889 compared with -7.089, while maintaining high solubility (0.857) and diversity (0.660).

Together with the SPSB2-targeted task, these results indicate that the permeability-enriched \mytool prior can be redirected to new protein structures through WP-DPO-driven optimization using only model-generated preference data, while maintaining favorable permeability and solubility.
This highlights the practical potential of PepALD for cyclic peptide design.

\begin{table*}[!b]
\centering
\footnotesize
\setlength{\tabcolsep}{0pt}
\begin{minipage}[t]{0.495\textwidth}
\centering
\textbf{(a) PepALD-generated cyclic peptides binding 6DN5}\\[1pt]
\renewcommand{\arraystretch}{1.08}
\begin{tabular*}{\linewidth}{@{\extracolsep{\fill}}l
S[table-format=-1.3]
S[table-format=1.3]
S[table-format=1.3]
S[table-format=-2.3]@{}}
\toprule
Model & \multicolumn{1}{c}{Permeability} & \multicolumn{1}{c}{Solubility} & \multicolumn{1}{c}{Diversity} & \multicolumn{1}{c}{Vina score} \\
\midrule
PepALD$_{\mathrm{perm}}$ & -5.804 & 0.794 & 0.656 & -5.832 \\
PepALD$_{\mathrm{SPSB2}}$ & -6.593 & 0.863 & 0.654 & -7.020 \\
Top $10\%$ & -6.622 & 0.819 & 0.608 & -8.037 \\
Top $25\%$ & -6.619 & 0.836 & 0.634 & -7.746 \\
Reference ligand & -7.181 & 0.785 & \multicolumn{1}{c}{--} & -7.627 \\
\bottomrule
\end{tabular*}
\end{minipage}\hfill%
\begin{minipage}[t]{0.495\textwidth}
\centering
\textbf{(b) PepALD-generated cyclic peptides binding 9BT3}\\[1pt]
\renewcommand{\arraystretch}{1.08}
\begin{tabular*}{\linewidth}{@{\extracolsep{\fill}}l
S[table-format=-1.3]
S[table-format=1.3]
S[table-format=1.3]
S[table-format=-2.3]@{}}
\toprule
Model & \multicolumn{1}{c}{Permeability} & \multicolumn{1}{c}{Solubility} & \multicolumn{1}{c}{Diversity} & \multicolumn{1}{c}{Vina score} \\
\midrule
PepALD$_{\mathrm{perm}}$ & -5.804 & 0.794 & 0.656 & -7.641 \\
PepALD$_{\mathrm{MtbCM}}$ & -5.889 & 0.857 & 0.660 & -8.645 \\
Top $10\%$ & -5.987 & 0.876 & 0.631 & -10.324 \\
Top $25\%$ & -5.934 & 0.865 & 0.631 & -9.781 \\
Reference ligand & -7.089 & 0.917 & \multicolumn{1}{c}{--} & -9.730 \\
\bottomrule
\end{tabular*}
\end{minipage}
\caption{Target-specific optimization metrics.
(a) Optimization of PepALD-generated cyclic peptides binding 6DN5.
(b) Optimization of PepALD-generated cyclic peptides binding 9BT3.
Each panel reports scalar metrics for PepALD$_{\mathrm{perm}}$, the final task policy, top-ranked generated subsets, and the reference ligand.
Permeability, solubility, and diversity are higher-is-better; Vina score is lower-is-better.}
\label{tab:spsb2_task}
\label{tab:mtbcm_task}
\label{tab:specific_task_metrics}
\end{table*}

\begin{figure*}[!t]
\centering
\small
\begingroup
\newcommand{\dockposepanel}[9]{%
\begin{minipage}[t]{0.315\textwidth}
\centering
\textbf{#1. #2}\\[3pt]
\begin{tikzpicture}
\node[inner sep=0pt] (poseimage)
{\includegraphics[width=\linewidth,height=1.25in,keepaspectratio,#4]{#3}};
\node[
anchor=south east,
fill=white,
fill opacity=0.90,
text opacity=1,
inner xsep=2.2pt,
inner ysep=1.6pt,
font=\tiny\sffamily,
align=left
] at ([xshift=-1pt,yshift=1pt]poseimage.south east)
{\textbf{Vina score}\\#5\\[-0.5pt]\textbf{Permeability}\\#6\\[-0.5pt]\textbf{Solubility}\\#7};
\end{tikzpicture}\\[-1pt]
\parbox{\linewidth}{\raggedright\tiny\ttfamily
\detokenize{#8}\\[-1pt]\detokenize{#9}\par}
\end{minipage}%
}
\begin{minipage}[t]{0.315\textwidth}
\centering
\textbf{a. SPSB2-targeted WP-DPO optimization}\\[3pt]
\includegraphics[width=\linewidth,height=1.25in,keepaspectratio]{result_data/case1_samples_vina/case1_vina_violin.pdf}
\end{minipage}\hfill
\dockposepanel{b}{SPSB2-CP1}
{result_data/3d_plot/case1/ligand4.png}
{trim=0 0 60 0,clip}
{-8.436}{-6.300}{0.851}
{PEPTIDE1{A.Mono84.Y.Me_Phe(3-Cl).dP.G.Mono84}}
{$PEPTIDE1,PEPTIDE1,1:R1-7:R2$$$}\hfill
\dockposepanel{c}{SPSB2-CP2}
{result_data/3d_plot/case1/ligand7.png}
{trim=0 0 60 0,clip}
{-8.256}{-6.576}{0.839}
{PEPTIDE1{Me_Bal.P.Mono2.T.P.Y.Phe(4-CF3)}}
{$PEPTIDE1,PEPTIDE1,1:R1-7:R2$$$}
\par\vspace{7pt}
\begin{minipage}[t]{0.315\textwidth}
\centering
\textbf{d. MtbCM-targeted WP-DPO optimization}\\[3pt]
\includegraphics[width=\linewidth,height=1.25in,keepaspectratio]{result_data/case2_samples_vina/case2_vina_violin.pdf}
\end{minipage}\hfill
\dockposepanel{e}{MtbCM-CP1}
{result_data/3d_plot/case2/ligand3.png}
{trim=410 0 360 0,clip}
{-11.067}{-5.936}{0.794}
{PEPTIDE1{Me_Bal.Mono84.P.L.meA.dF.D}}
{$PEPTIDE1,PEPTIDE1,1:R1-7:R2$$$}\hfill
\dockposepanel{f}{MtbCM-CP2}
{result_data/3d_plot/case2/ligand4.png}
{trim=410 0 360 0,clip}
{-10.034}{-5.617}{0.834}
{PEPTIDE1{Mono6.T.Sar.meF.P.G.D}}
{$PEPTIDE1,PEPTIDE1,1:R1-7:R2$$$}
\caption{Target-specific WP-DPO optimization redirects PepALD toward receptor-favorable macrocyclic peptides while retaining permeability-related properties.
(a) Vina score distributions for SPSB2-targeted cyclic peptide generation, targeting the SPSB2--iNOS interface represented by 6DN5, from PepALD$_{\mathrm{perm}}$ through iterative WP-DPO rounds.
(b,c) Representative optimized SPSB2 cyclic-peptide candidates docked to 6DN5.
(d) Vina score distributions for MtbCM-targeted cyclic peptide generation, targeting Mtb chorismate mutase (MtbCM) represented by 9BT3, under the same optimization protocol.
(e,f) Representative optimized MtbCM cyclic-peptide candidates docked to 9BT3.}
\label{fig:specific_task_optimization}
\endgroup
\end{figure*}
